\chapterrule
\section{预训练已经学会续写，为什么还能临时学会新任务？}
\label{intro:23aa-pretraining-icl}
\chaptercompass{新任务能否不更新权重就学会格式？}{把任务说明、示例和新输入放入同一段可见上下文。}{上下文 Token $\to$ 下一 Token 概率}{ICL 保持 $\theta$ 不变；SFT 才使用梯度调整参数}
\modelpipelinefocus{8}{上下文任务与 ICL}
\phantomsection\label{concept:pretraining}
\phantomsection\label{concept:in-context-learning}


上一部分用 CNN 与 ViT 说明同一套向量和注意力方法如何处理图像。
现在进入\textbf{助手面对当前请求}的场景：
不进行梯度下降时，模型为什么仍能根据上下文示例模仿任务格式？

\subsection{预训练扩大的是统计经验，不是写死的知识清单}

把许多文档转成 Token 序列，使用因果 Mask，
每个位置尝试预测后续 Token，再对大量位置和样本的损失做优化，
这就是自回归\textbf{预训练}的一种基本训练目标：
\[
L_{\rm pretrain}(\theta)=
-\frac1N\sum_{(x,y)}\log p_\theta(y\mid x).
\]
$x$ 表示位置之前的可见文本，$y$ 表示下一 Token；
$N$ 是参与计分的训练位置总数，并非文档数。
它和本书的三位置交叉熵\textbf{在目标形式上相同}；
区别在数据范围、训练预算与模型容量，而不是突然改成另一种数学魔法。

大规模训练会调整 Embedding、Attention、FFN 等参数，
使模型对训练分布中的多种语言模式形成统计能力。
它可以帮助模型迁移到新文本，但\textbf{并不保证事实正确、持续更新或理解每一条规则}。

\subsection{训练资料从哪里来：来源、清洗、去重与配比}
\phantomsection\label{concept:pretraining-data}
预训练前的关键工作不是先选网络层数，而是把可使用的资料变成
\textbf{有来源、有边界的 Token 训练样本}。
例如原始语料可能混有网页正文、开放许可书籍、代码与问答；
它们的格式、质量、许可和语言比例都不同。

\begin{enumerate}
\item \textbf{来源与授权：}记录来源、许可和使用范围，
剔除不适合训练的个人敏感信息与不允许使用的数据；
\item \textbf{数据清洗：}去掉破损编码、导航栏、广告模板、
无意义重复行，保留有用的正文与必要上下文；
\item \textbf{精确及近似去重：}通过内容哈希删除完全重复文档，
再用分片相似度等方法发现大量改写或镜像转载的近重复文本；
\item \textbf{质量与配比：}平衡语言、领域、代码和教程等来源，
避免高频模板挤占训练信号，并持续测量各来源的贡献；
\item \textbf{切分与污染排查：}构造训练／验证／测试边界，
核对基准测试题和答案是否已混入训练数据。
\end{enumerate}

\localfixture{一组可数清的语料筛选例}
假定收集 $100$ 篇文档，其中有 $10$ 篇精确重复、
$5$ 篇近重复、$5$ 篇不符合使用条件，且这些集合互不重叠。
清理后为 $100-10-5-5=80$ 篇。若再从中留出 $8$ 篇验证，
参与本轮训练的为 $72$ 篇；这只是文档数，
并不等于 Token 数或梯度更新步数。
真实去重集合可能重叠，应先去重合并后再计数。

这种数据治理决定模型接触哪些表达与事实，也决定偏差与污染风险；
注意力机制本身不会替训练过程筛选可靠来源。

在第~\ref*{intro:22b-token-to-output}~章，
“春天 花 开了 。”只是一段下一 Token 训练样本；
这里把\textbf{句子续写}与\textbf{今天某处是否真的开花}分开：
前者检验模型预测，后者需要当日可核验的外部观察。

\subsection{零样本、少样本：能在不改权重时给出示例吗？}
\phantomsection\label{concept:few-shot}

以接下来将贯穿系统章节的\textbf{季度运营报告}任务为例。
先提出一个纯文本格式要求：“阅读指标，输出季度与收入，附原文出处”；
这是\textbf{零样本提示}（Zero-shot）的一种形式。
再在 Prompt 中给出两条例子：

\begin{quote}
输入：一季度收入为 100 百万元。\\
输出：季度=Q1；收入=100；单位=百万元。\\
输入：二季度收入为 120 百万元。\\
输出：季度=Q2；收入=120；单位=百万元。
\end{quote}

然后输入“第三季度收入为 132 百万元”，
希望模型按示例输出相同格式。
这叫\textbf{上下文学习}（In-context Learning, ICL）中
常见的 Few-shot（少样本）使用方式。
它\textbf{不是把这两条文本变成梯度去更新网络}：
预测时模型参数 $\theta$ 保持不变，
变化的是作为条件的 Prompt：
\[
p_\theta(y\mid\underbrace{\text{任务说明、示例、新输入}}_{\text{上下文}}).
\]
模型可以模仿格式，也可能漏掉单位或编造不存在的来源；
示例只是条件，不是正确性的保证。

\subsection{一条清楚的边界：上下文 vs 权重 vs 外部资料}

\begin{itemize}
\item \textbf{预训练/微调}：计算损失和梯度，改变权重 $\theta$；
持久改变模型参数，但更新成本与风险高。
\item \textbf{ICL/Prompt}：只改变本次上下文，不改变 $\theta$；
会受到上下文长度、顺序、示例质量影响。
\item \textbf{RAG}：从系统允许的数据中检索片段并放入上下文；
它也通常不改模型权重，来源与准确性仍须核验。
\end{itemize}
尤其要防止误解：从几条新示例模仿格式，不代表模型在本次推理中
运行了与预训练相同的参数优化步骤。

\subsection{为什么大规模预训练也不是万能钥匙？}

例如训练文本大多只讨论到某一日期，
模型参数里没有自动更新的未来报表。
即使见过财务术语，也不能凭语言统计
代替\textbf{当前文件读取、来源核对与准确计算}。
Few-shot 能让某些任务在不更新参数的情况下改善表现，
却不是“给一个例子就必然学会”的定律。
经典研究可参见 Brown 等（2020），
\href{https://arxiv.org/abs/2005.14165}{Language Models are Few-Shot Learners}；
其中的实验结论有其模型、数据和评测范围。

\chapterchecktasks{23aa-pretraining-icl}{正文基础三题}{
\item \textbf{比较：}预训练、SFT、Few-shot ICL 中哪些需要修改参数？“有训练目标”与“在提示里临时提供例子”的差别在哪里？
\item \textbf{手算：}Prompt 中放两条示例及一条待回答输入，模型本次能看到几条任务记录？若另开一个没有带入示例的新会话，原来的两条例子是否仍在上下文？
\item \textbf{找错：}某人把新的 Q3 数字复制进 Few-shot 示例后宣称“模型权重已经永久更新”。指出错误，并说出如何得到可信 Q3 资料。
}
\chapterrecallanswer{ICL 改变一次请求的上下文，SFT 才通过训练更新参数；真实的新资料还需要合法的检索或输入。}



\begin{conclusion}
预训练用海量预测题改变权重；ICL 利用已训练模型与当前上下文，
通常不改变权重。下一道缺口：
即使模型能模仿任务格式，我们又怎样让它更稳定地遵守指令与评价标准？
\end{conclusion}
